\begin{frame}[plain]
\vfill{\usebeamerfont{title}\color{AlgoBlue}선택과 순서통계량\par}
\vspace{1mm}{\Large Selection and Order Statistics\par}
\vspace{3mm}{\color{AlgoOrange}\rule{38mm}{1.2pt}\par}
\vspace{5mm}{\large Quickselect에서 Median of Medians까지\par}
\vspace{8mm}{\small Algorithm Lecture Notes\par}\vfill
\end{frame}
\begin{frame}{학습 목표}
\begin{itemize}
  \item $i$th order statistic을 정의하고 subarray-relative rank와 array index를 구분한다.
  \item Quickselect를 추적하며 오른쪽 재귀의 rank를 정확히 갱신한다.
  \item partition property와 recursive invariant로 Quickselect의 total correctness를 설명한다.
  \item fixed pivot의 worst case와 randomized pivot의 expected behavior를 구분한다.
  \item Median of Medians의 pivot 보장과 BFPRT recurrence를 설명하고 전략을 선택한다.
\end{itemize}
\end{frame}
\begin{frame}{왜 전체 정렬이 과할 수 있는가?}
\begin{columns}[T]\column{.47\textwidth}\begin{block}{질문 하나}중앙값 하나, 상위 10\% 경계, $k$번째 값만 필요하다.\end{block}\column{.47\textwidth}\begin{alertblock}{정렬의 비용}전체 $n$개의 상대 순서를 모두 만든다: $\Theta(n\log n)$.\end{alertblock}\end{columns}
\selecttakeaway{Selection은 필요한 rank 하나를 찾고 불필요한 순서는 만들지 않는다.}
\end{frame}
\begin{frame}{Order Statistic은 어디에 쓰이는가?}
\begin{itemize}\item latency의 95th percentile\item 성적 상위 10\%의 경계\item 영상 filter의 지역 median\item 대규모 데이터의 $k$th smallest\item streaming에서는 exact selection 대신 sketch가 필요할 수도 있음\end{itemize}
\end{frame}
\begin{frame}{Lecture 4 Road Map}
\centering\begin{tikzpicture}[every node/.style={tree node,align=center,font=\small,minimum width=25mm}]
\node[chosen](s)at(0,0){Selection};
\node(st)at(-4.4,-1.2){Sort then Select};
\node(q)at(0,-1.2){Partition-based\\Selection};
\node(d)at(4.4,-1.2){DeterministicSelect};
\node(f)at(-1.7,-2.7){Fixed-pivot\\Quickselect};
\node(r)at(1.7,-2.7){RandomizedSelect};
\node(m)at(4.4,-2.7){Median of Medians\\worst $\Theta(n)$};
\draw[-Latex](s)--(st);\draw[-Latex](s)--(q);\draw[-Latex](s)--(d);
\draw[-Latex](q)--(f);\draw[-Latex](q)--(r);\draw[-Latex](d)--(m);
\end{tikzpicture}
\vfill\muted{Quickselect는 한쪽만 남기는 framework이고, 시간 보장은 pivot policy에 따라 달라진다.}
\end{frame}
